% language=us

\environment context-2022-style

\startdocument
  [title={Injecting stuff},
   banner={fundamental or not},
   location={context\enspace {\bf 2022}\enspace meeting}]

\starttitle[title=Where to draw a line]

What can be done in \LUA\ should be done in \LUA, unless:

\startitemize
    \startitem
        There is a too tight interaction with the builders: par, page, math.
    \stopitem
    \startitem
        Information is only available at specific stages and is hard to
        communicate.
    \stopitem
    \startitem
        The callback overhead would be too large due to either too many
        slightly different places to call and/or contexts to check.
    \stopitem
    \startitem
        Something really should be part of the core engine (also at the
        macro, input, tracing, etc level).
    \stopitem
\stopitemize

\blank[3*big]

{\em This is per 2022, so expect things to have matured size then. The low-level
manuals get updated as we go forward and more examples are added to the test
suite. The mechanisms mentioned here are mostly feature complete in 2026.}

\stoptitle

\starttitle[title=Boxes]

\startitemize

\startitem
    Traditionally \TEX\ has two kind of boxes and they can be set, copied,
    flushed either or not unpacked.
\stopitem

\startitem
    Most box processing is influenced by (optional) callbacks: the most
    significant one is the one that deals with fonts.
\stopitem

\startitem
    These callbacks can be avoided by packing instead of boxing (the full
    package).
\stopitem

\startitem
    Repacking can be avoided by using some manipulators (a chicken|-|egg
    situation).
\stopitem

\blank[3*big]

\starttyping
lowlevel-boxes
\stoptyping

\stopitemize

\stoptitle

\starttitle[title=Leaders]

\startitemize

\startitem
    Leaders are very \TEX ie. They are bound to glue and have basic content like
    boxes and rules, although in \LUAMETATEX\ we can also have glyphs and math.
\stopitem

\startitem
    They fill the space calculated for the glue by repeated content.
\stopitem

\stopitemize

\blank[3*big]

\starttyping
lowlevel-rules   % to be done
lowlevel-leaders % to be done
\stoptyping

\stoptitle

\starttitle[title=Inserts]

\startitemize

\startitem
    These are content blobs that travel around with the rest. Think footnotes.
\stopitem

\startitem
    They are part of vertical calculations and therefore vertical splitting like
    page breaks.
\stopitem

\startitem
    They have side effects and complicate for instance faking multi column
    processing.
\stopitem

\startitem
    The traditional engine (naturally) stores insert related properties in
    registers. In \LUAMETATEX\ we have a dedicated name space.
\stopitem

\startitem
    We move around inserts in a different way. This avoids them being not seen as
    in regular \TEX\ when they are buried deeply.
\stopitem

\stopitemize

\blank[3*big]

\starttyping
lowlevel-inserts
\stoptyping

\stoptitle

\starttitle[title=Marks]

\startitemize

\startitem
    Marks are special nodes that carry token lists around in the text. Think
    section titles in headers and footers.
\stopitem

\startitem
    They are invisible but can be accessed in the output routine where
    they have been set according the split page box.
\stopitem

\startitem
    As with inserts in \LUAMETATEX\ they travel around a bit different. They also
    bubble up.
\stopitem

\startitem
    In \MKII\ we always used marks extensively for state info (\ETEX\ made
    that more efficient). In \MKIV\ we emulated what is now more natural in
    \LUAMETATEX\ and \MKIX\ (\LMTX).
\stopitem

\stopitemize

\blank[3*big]

\starttyping
lowlevel-marks
\stoptyping

\stoptitle

\starttitle[title=Vadjust]

\startitemize

\startitem
    This mechanism permits pushing stuff in front of or after the current
    line.
\stopitem

\startitem
    It is hard to use because it interferes badly with spacing.
\stopitem

\startitem
    As side effect of some math experiments this mechanism got extended
    with better spacing control.
\stopitem

\startitem
    It might evolve a bit more (which is no problem because we use it low level
    and we can adapt).
\stopitem

\stopitemize

\blank[3*big]

\starttyping
lowlevel-adjust % todo
vadjust-001.tex
\stoptyping

\stoptitle

\starttitle[title=Local boxes]

\startitemize

\startitem
    This difficult and confusing feature comes from \ALEPH\ and can bind
    content to the left and right end of lines.
\stopitem

\startitem
    The dimensions are taken into account at line breaks.
\stopitem

\startitem
    It's semi global nature makes it hard to deal with.
\stopitem

\startitem
    They have been redone in \LUAMETATEX\ and we now use them for some
    experiments.
\stopitem

\stopitemize

\blank[3*big]

\starttyping
lowlevel-localboxes
\stoptyping

\stoptitle

\starttitle[title=Anchors]

\startitemize

\startitem
    This is a new mechanism that can be seen as side track of the math
    engine overhaul.
\stopitem

\startitem
    It is a way to bind boxes together and is something that the backend
    has to deal with.
\stopitem

\startitem
    So it is more about carrying around information. It could have been
    attributes but this is more efficient when used large scale.
\stopitem

\startitem
    Usage is not standardized.
\stopitem

\stopitemize

\blank[3*big]

\starttyping
ontarget-anchoring
anchoring-001.tex
\stoptyping

\stoptitle

% \starttitle[title=Paragraphs]
% \stoptitle

% \starttitle[title=Groups]
% \stoptitle

\stopdocument
